\begin{table*}[ht!]
\small
\caption{Examples of model predictions. Gold leaf is highlighted in \textbf{bold} in Table of Contents}
\label{table:pred_examples}
\resizebox{\textwidth}{!}{%
\begin{tabular}{l|l}
\multicolumn{2}{c}{\textbf{Example 1:}} \\
\toprule
\multicolumn{2}{l}{\textbf{Book}: \href{https://open.umn.edu/opentextbooks/textbooks/1621}{The Whole Child: Development in the Early Years}}\\
\multicolumn{2}{l}{\textbf{Query}: How do preschoolers react when caregivers and teachers belittle their autonomous actions?}\\
\midrule
\textbf{Table of Contents} & \textbf{Content}\\
$\cdots$\\
2 Chapter Two: Theorists and Theories of Development & Initiative vs. Guilt (Preschool Years)\\
leaf 2.1 Theories of Development & The development of courage and independence ...\\
2.2 Psychosocial Theory & ...\\
leaf 2.2.1 Trust vs. Mistrust (Infancy)& ...\\
leaf 2.2.2 Autonomy vs. Shame/Doubt (Toddlerhood) &...\\
\textbf{leaf 2.2.3 Initiative vs. Guilt (Preschool Years)}\\
leaf 2.2.4 Industry vs. Inferiority (Elementary Years) & If caregivers and preschool teachers encourage and support\\
leaf 2.2.5 Identity vs. Role-Confusion (Adolescence) & children’s efforts while also helping them make realistic \\
leaf 2.2.6 Intimacy vs. Isolation (Early Adulthood) & and appropriate choices, children develop a healthy sense of initiative in \\
leaf 2.2.7 Generativity vs. Stagnation (Middle Adulthood)& planning and undertaking activities If, instead, adults
discourage the\\
leaf 2.2.8 Integrity vs. Despair (Older Adulthood) & pursuit of independent activities or dismiss them as silly and \\
$\cdots$ & bothersome, children develop guilt about their needs and desires.\\
3 Chapter Three: Domains in Development\\
\midrule
\multicolumn{2}{l}{\textbf{Mistral}: 8.3.3 Moral Development, leaf 7.3.2 Hitting/Scratching and Temper Tantrums, \textbf{leaf 2.2.3 Initiative vs. Guilt (Preschool Years)}}\\
\multicolumn{2}{l}{\textbf{BM25}: 8.3.2 Moral Development}\\
\multicolumn{2}{l}{\textbf{DPR: }2.2.2 Autonomy vs. Shame/Doubt (Toddlerhood)}\\
\multicolumn{2}{l}{\textbf{DSI}: 2.2.2 Autonomy vs. Shame/Doubt (Preschool Years)}\\
\multicolumn{2}{l}{\textbf{{\sys}}: \textbf{2.2.3 Initiative vs. Guilt (Preschool Years)}}\\
\midrule
\bottomrule
\multicolumn{2}{c}{\textbf{Example 2:}} \\
\toprule
\multicolumn{2}{l}{\textbf{Book}: \href{https://open.umn.edu/opentextbooks/textbooks/1013}{Nursing Fundamentals}}\\
\multicolumn{2}{l}{\textbf{Query}: In what way does a chronic illness affect an elderly person's ability to perform daily activities?}\\
\midrule
\textbf{Table of Contents} & \textbf{Content}\\
$\cdots$\\
leaf 30.5 Spiritual Care of Self & Applying the Nursing Process\\
leaf 30.6 Putting It All Together& ...\\
31 Care of the Older Adult & It is also important to consider the impact of chronic disease\\
leaf 31.1 Care of the Older Adult Introduction & on
their ability to function and complete Activities of Daily Living\\
leaf 31.2 Basic Concepts &(ADLs). older adults who are able to perform ADLs without assistance\\
leaf \textbf{31.3 Applying the Nursing Process} & consider themselves healthy.\\
\midrule
\multicolumn{2}{l}{\textbf{Mistral}: \textbf{31.3 Applying the Nursing Process}}\\
\multicolumn{2}{l}{\textbf{BM25}: 11.2 Sensory Impairments Basic Concepts}\\
\multicolumn{2}{l}{\textbf{DPR: } 31.2 Basic Concepts}\\
\multicolumn{2}{l}{\textbf{DSI}: 31.2 Basic Concepts}\\
\multicolumn{2}{l}{\textbf{{\sys}}: \textbf{31.3 Applying the Nursing Process}}\\
\midrule
\bottomrule
\multicolumn{2}{c}{\textbf{Example 3:}} \\
\toprule
\multicolumn{2}{l}{\textbf{Book}: \href{https://open.umn.edu/opentextbooks/textbooks/1435}{Introduction to Comparative Government and Politics}}\\
\multicolumn{2}{l}{\textbf{Query}: What are irregular armed organizations and how are they used by states?}\\
\midrule
\textbf{Table of Contents} & \textbf{Content}\\
9 Chapter 5: Non-Democracies and Democratic Backsliding\\
$\cdots$ & Another powerful instrument of repression are paramilitaries.\\
9.2 Strategies for staying in power & These refer to groups with access 
to military-grade weapons and training\\
leaf 9.2.1 \textbf{Institutional channels}& yet they are not part of the national military.\\
leaf 9.2.2 Cultural and ideological controls & They are “irregular armed organizations that carry out acts of violence\\
9.3 Varieties of non-democracy & against civilians on behalf of a state,” \\
$\cdots$\\
\midrule
\multicolumn{2}{l}{\textbf{Mistral}: 9.3.5 Illiberal and hybrid regimes}\\
\multicolumn{2}{l}{\textbf{BM25}: 15.3.1 Insurgencies/Civil Wars}\\
\multicolumn{2}{l}{\textbf{DPR}: 15.2.2 External State-Sponsored Political Violence (State-Sponsored Terrorism)}\\
\multicolumn{2}{l}{\textbf{DSI}: 15.2.1 Internal State-Sponsored Political Violence (Government Terrorism)}\\
\multicolumn{2}{l}{\textbf{{\sys}}: \textbf{9.2.1 Institutional channels}}\\
\end{tabular}
}
\end{table*}

